\chapter{A Bigger World: Arrays, Structs, and Tile Maps}
\label{ch:c2}

The Pong of Chapter~\ref{ch:c1} held a handful of separate variables. Real games
hold many things at once: rows of bricks, swarms of enemies, a whole landscape of
ground and platforms. This chapter introduces the two C features that make such
games manageable: the \emph{array}, which holds many values of the same kind, and
the \emph{struct}, which groups related values into one named thing, and applies
them to a tile-based platform game built on the platform's own example
\citep{prg32repo,prg32tutorialc}.

\section{Arrays: many values under one name}

An array is a numbered row of values of the same type. Where Pong needed one
paddle position, Breakout needs a position, or simply a present/absent flag, for
each of many bricks, and writing a separate variable for each would be absurd. An
array lets you store them together and reach any one by its index.

\lstinputlisting[
    style=c,
    caption={%\fname{brick\_flags.c} \protect\
       An array of brick flags. Index $i$ holds whether brick $i$ is still present.}
]{examples/ch08/brick_flags.c}

\begin{prgconcept}
An array stores many values of one type in a numbered sequence; \code{bricks[i]}
is the value at index \code{i}, counting from zero. A \code{for} loop visits every
element in turn. Together, the array and the loop replace dozens of separate
variables and the repetitive code that would manage them.
\end{prgconcept}

The framework itself uses small arrays to describe artwork. An 8-by-8 tile: a
patch of ground, a coin or a hazard, is stored as an array of eight bytes, one per
row, where each bit is a pixel. The platformer example defines its tiles exactly
this way \citep{prg32repo}.

\lstinputlisting[
    style=c,
    caption={%\fname{coin\_tile.c} \protect\
       An 8-by-8 tile as an array of eight bytes. Each byte is one row; each bit is one pixel. This is the platformer's coin tile.}
]{examples/ch08/coin_tile.c}

\begin{prgaside}
Reading the bits of \code{tile\_coin} as ones and zeros, top to bottom, sketches a
small round coin: the rows are narrow at top and bottom and wide in the middle.
Defining graphics as bit patterns in arrays is exactly how the home computers of
the retro era stored their character sets and sprites, with artists working in
hexadecimal on grid paper. The technique is old, compact, and still perfectly
legible.
\end{prgaside}

\section{Structs: grouping what belongs together}

A moving character has several attributes that always travel together: its
position, its velocity, its size, its state. Keeping them in separate arrays that
must be kept in step is error-prone. A \emph{struct} bundles them into a single
named value, so the whole character is one thing you can pass around. The framework
defines such a struct for platform-game characters, \code{prg32\_platform\_actor\_t}
\citep{prg32repo}.

\lstinputlisting[
    style=c,
    caption={%\fname{platform\_actor\_struct.c} \protect\
       The framework's actor struct (from \texttt{prg32.h}). Position, velocity, size, and state, grouped into one value.}
]{examples/ch08/platform_actor_struct.c}

\begin{prgconcept}
A struct groups related values into one named type. With
\code{prg32\_platform\_actor\_t player;} you create a single \code{player} whose
position is \code{player.x}, velocity \code{player.vx}, and so on. Passing one
struct is cleaner and safer than passing six loose variables, and it scales: an
array of structs is a crowd of characters.
\end{prgconcept}

\section{A tile map: the world as a grid}

A platform game's world is too large to draw as individual rectangles. Instead it
is a \emph{tile map}: a grid in which each cell names a tile to draw there. The
framework provides a two-layer playfield for this, and you fill it by placing tile
identifiers at grid coordinates with \code{prg32\_playfield\_put} \citep{prg32repo}.
The platformer example builds its landscape with small helpers that draw runs and
columns of tiles: ordinary loops over the grid.

\lstinputlisting[
    style=c,
    caption={%\fname{put\_run.c} \protect\
       Filling a run of ground tiles with a loop. The platformer uses helpers like this to lay out its world.}
]{examples/ch08/put_run.c}

\begin{prgconcept}
A tile map represents a large world compactly: instead of remembering every pixel,
it remembers, for each grid cell, which small reusable tile belongs there. Drawing
the world becomes drawing the grid; building the world becomes filling the grid
with loops. This is how side-scrolling games fit sprawling levels into tiny
memory.
\end{prgconcept}

\section{Behaviour from tile flags}

A tile is not only a picture; it can carry \emph{behaviour}. The framework lets you
attach flags to a tile identifier, solid, platform, hazard, collectable, and then
the physics helpers know how a character should interact with it. You declare, for
instance, that tile~2 is solid, and from then on a character cannot walk through
it \citep{prg32repo}.

\lstinputlisting[
    style=c,
    caption={%\fname{tile\_flags.c} \protect\
       Giving tiles behaviour. Marking a tile solid makes the physics helper treat it as ground or wall.}
]{examples/ch08/tile_flags.c}

\section{Putting it together: an actor in a world}

With an actor struct, a tile map, and tile flags, a platform game's per-frame
logic becomes remarkably short. The framework's \code{prg32\_platform\_actor\_step}
reads the input, applies gravity, moves the actor, and resolves collisions against
the solid tiles, updating the struct's \code{state} flags so your code can react
to what happened \citep{prg32repo,prg32tutorialc}.

\lstinputlisting[
    style=c,
    caption={%\fname{platform\_physics.c} \protect\
       One frame of platform physics. The helper does the movement and collision; your code reads the resulting state.}
]{examples/ch08/platform_physics.c}

Notice the economy. A single call advances the whole character: the joystick moves
it, gravity pulls it down, the solid tiles stop it, and the returned \code{state}
tells you if it touched a hazard. The camera-follow call then scrolls the tile map
so the character stays in view. The same behaviour written from first principles
would be a page of code; the struct and the helper compress it because they capture
exactly the data and the operation the problem needs.

\begin{prgwarn}
Pass the \emph{address} of the actor, \code{\&player}, to the stepping and
follow helpers, not the struct by value. These functions must \emph{change} the
actor, its position and state, and in C a function can only change a value the
caller can see if it is given that value's address. Forgetting the \code{\&} is a
common error; the symptom is an actor that never moves, because the function
updated a private copy.
\end{prgwarn}

\begin{prgtargets}
\textbf{Both targets.} The platformer's tile map, gravity, and camera all render on
the QEMU virtual screen, so you can build and tune the whole level on a laptop. On
the ESP32\nobreakdash-C6 you then play it with the joystick and hear the hazard beep
\citep{prg32qemu}. The struct, the array of tiles, and the physics call are
identical across the two targets.
\end{prgtargets}

\begin{prgtry}
Open \fname{examples/games/platformer/c/game.c} and find where it defines the coin
tile and where it marks a tile collectable. Add one new platform to the world by
calling \code{put\_run} with new coordinates, rebuild for QEMU, and confirm your
platform appears. Then change the \code{gravity} argument of
\code{prg32\_platform\_actor\_step} from \code{1} to \code{2} and describe how the
jump feels different. One change, observed and recorded \citep{prg32repo}.
\end{prgtry}

\begin{prgcheck}
You can now hold many values in arrays, group related values in structs, represent
a world as a tile map filled by loops, and give tiles behaviour with flags. These
are the data-structuring ideas a programming course exists to teach, learned here
on a game you can see and play. Chapter~\ref{ch:ccap} brings them together into a
capstone you extend and ship as a cartridge.
\end{prgcheck}
